% =====================================================================
% PARTE II — I tipi
% =====================================================================
\part{I tipi}

\noindent\itshape Otto capitoli, sempre F-I-E-C.
$\mathbb{B}$ (\S\ref{cap:bool}) per vedere lo schema nel caso pi\`u puro, senza
ricorsione. $\mathbb{N}$ (\S\ref{cap:nat}) per il salto all'induzione vera.
$\Pi$ (\S\ref{cap:pi}) per la decisione fondativa di MLTT, e $\Pi$ in azione
(\S\ref{cap:pi-azione}) per le due passate che la sintassi nasconde. $\forall$
(\S\ref{cap:forall}) come zucchero. $\Sigma$ (\S\ref{cap:sigma}) come duale.
Il tipo identit\`a $\_\identita\_$ (\S\ref{cap:identita}) come sezione cardine
--- l\`i si chiude la distinzione fra le due uguaglianze. $\mathrm{Vec}\,A\,n$
(\S\ref{cap:vec}) come capstone, dove $\mathbb{N}$, $\Pi$ e la dipendenza si
incastrano per la prima volta su un valore numerico.\upshape

\bigskip

% ---------------------------------------------------------------------
\chapter{$\mathbb{B}$ --- il caso pi\`u puro}\label{cap:bool}

Cominciamo da $\mathbb{B}$. \`E il tipo induttivo pi\`u piccolo: due costruttori, zero argomenti, zero ricorsione. L'eliminazione qui \`e \textbf{puro case-split} --- non c'\`e ipotesi induttiva, perch\'e non c'\`e ricorsione su cui ricorrere. Vederlo prima di $\mathbb{N}$ rende esplicito che la ricorsione di $\mathbb{N}$ \`e qualcosa \emph{in pi\`u}, non parte intrinseca di ``eliminare''.

\section*{F --- Formazione}
\[
\inferrule{ }{\mathbb{B} : \Type}
\]
Una regola senza premesse: $\mathbb{B}$ esiste, punto.

\section*{I --- Introduzione}
\[
\inferrule{ }{\tto : \mathbb{B}}
\qquad
\inferrule{ }{\ff : \mathbb{B}}
\]
Due costruttori, entrambi senza premesse.

\section*{E --- Eliminazione (metalinguaggio)}
\[
\inferrule{
b : \mathbb{B} \\
C : \mathbb{B} \to \Type \\
c_t : C\,\tto \\
c_f : C\,\ff
}{
\indB(b, C, c_t, c_f) : C\,b
}
\]

Letto a voce: ``per costruire qualcosa che dipende da un $b : \mathbb{B}$ generico --- cio\`e un abitante di $C\,b$ --- basta darmi il pezzo per $b = \tto$ e il pezzo per $b = \ff$; io metto tutto insieme''. $C$ si chiama \textbf{motivo} dell'eliminazione.

% NUOVO: paragrafo sul motivo, valido per tutti gli eliminatori
\begin{oss}[Il motivo $C$, una volta per tutti gli eliminatori]
Il motivo \`e il concetto pi\`u astratto della Parte~II, e ritorner\`a identico in $\indN$, $\indS$, $\Jelim$, $\indV$. Vale la pena fissarlo subito.

$C$ \`e \textbf{ci\`o che si vuole concludere}, espresso come funzione dal valore eliminato. Scriverlo come $C : \mathbb{B} \to \Type$ significa: ``per ogni $b$ possibile dir\`o cosa intendo costruire''. L'eliminatore promette di consegnare $C\,b$ per il $b$ effettivo; in cambio chiede un pezzo per ogni costruttore.

Due regimi tipici:
\begin{itemize}
\item \textbf{Motivo costante} ($C\,b \definitorio X$ per ogni $b$): si sta solo facendo una funzione $\mathbb{B} \to X$, niente di dipendente. L'eliminatore degenera a un ``if-then-else'' usuale.
\item \textbf{Motivo davvero dipendente} ($C\,b$ varia con $b$): si sta dimostrando una propriet\`a o costruendo un dato il cui \emph{tipo} cambia con il valore. Qui la dipendenza Cero~\ref{cero:dipendenza} viene effettivamente sfruttata.
\end{itemize}
Tutti gli eliminatori che seguono hanno la stessa anatomia: motivo + un caso per costruttore (pi\`u, dove c'\`e ricorsione, un'ipotesi induttiva).
\end{oss}

\section*{C --- Computazione (definitorie, Cero~\ref{cero:comp})}
\begin{align*}
\indB(\tto, C, c_t, c_f) &\definitorio c_t \\
\indB(\ff,  C, c_t, c_f) &\definitorio c_f
\end{align*}

Quando l'argomento \`e un costruttore concreto, l'eliminatore riduce al ramo corrispondente. Nient'altro.

\begin{oss}[Punto da fissare]
Confronta con quanto vedremo per $\mathbb{N}$ in \S\ref{cap:nat}: l\`i $c_t, c_f$ saranno sostituiti da un caso base e un \emph{passo} che riceve un'ipotesi induttiva. Qui no, perch\'e $\tto$ e $\ff$ non contengono altre $\mathbb{B}$. Eliminare = sempre case-split su tutti i costruttori. Indurre = case-split + ricorrere quando il costruttore contiene istanze del tipo stesso. $\mathbb{B}$ ha la prima cosa, non la seconda.
\end{oss}

% ---------------------------------------------------------------------
\chapter{$\mathbb{N}$ --- l'eliminazione diventa induzione}\label{cap:nat}

Stesso schema, ma ora un costruttore (\texttt{suc}) prende un argomento di tipo $\mathbb{N}$. Questo cambia l'eliminazione: nel ramo \texttt{suc k}, oltre a \texttt{k} stesso, riceverai anche $C\,k$ --- il risultato dell'eliminazione \emph{gi\`a calcolato} su \texttt{k}. \`E l'\textbf{ipotesi induttiva}.

\section*{F --- Formazione}
\[
\inferrule{ }{\mathbb{N} : \Type}
\]

\section*{I --- Introduzione}
\[
\inferrule{ }{\zero : \mathbb{N}}
\qquad
\inferrule{n : \mathbb{N}}{\suc\,n : \mathbb{N}}
\]

Due costruttori, ma uno (\texttt{suc}) ha una premessa: prende una $\mathbb{N}$ per produrne un'altra. \`E \textbf{questa} premessa che fa la differenza con $\mathbb{B}$.

\section*{E --- Eliminazione}
\[
\inferrule{
n : \mathbb{N} \\
C : \mathbb{N} \to \Type \\
b : C\,\zero \\
s : (k : \mathbb{N}) \to C\,k \to C\,(\suc\,k)
}{
\indN(n, C, b, s) : C\,n
}
\]

Confronto riga per riga con $\indB$:
\begin{itemize}
% MODIFICATO: uniformato \texttt{tt} -> \tto per coerenza con cap. Bool
\item per $\zero$ serve un valore $b : C\,\zero$ (come $c_t$ per $\tto$);
\item per $\suc\,k$ serve $s : (k : \mathbb{N}) \to C\,k \to C\,(\suc\,k)$ --- non un valore, ma una funzione che prende $k$ \emph{e l'ipotesi induttiva $C\,k$} e produce $C\,(\suc\,k)$.
\end{itemize}

Il $C\,k$ in input a $s$ \`e dove l'eliminazione di $\mathbb{N}$ diventa \textbf{ricorsione primitiva}. \`E letteralmente l'ipotesi induttiva.

\section*{C --- Computazione}
% MODIFICATO: spaziatura uniforme negli argomenti di \indN
\begin{align*}
\indN(\zero,\,C, b, s)    &\definitorio b \\
\indN(\suc\,k,\,C, b, s)  &\definitorio s\,k\,(\indN(k, C, b, s))
\end{align*}

Il caso $\suc\,k$ \`e dove appare la \textbf{ricorsione}: $\indN$ chiama s\'e stesso su $k$, strutturalmente pi\`u piccolo di $\suc\,k$. Questo garantisce terminazione.

% NUOVO: indN in azione --- la somma come esempio concreto di ipotesi induttiva
\section*{$\indN$ in azione --- la somma}

L'addizione \`e l'esempio canonico per vedere l'ipotesi induttiva al lavoro. Definendola tramite $\indN$ con motivo costante $C \definitorio (\lambda\,\_\,.\,\mathbb{N} \to \mathbb{N})$:
\[
n + m \;\definitorio\; \indN\!\bigl(n,\; \lambda\,\_\,.\,\mathbb{N},\; m,\; \lambda\,k\,\mathit{rec}.\,\suc\,\mathit{rec}\bigr)
\]
Letta a voce: ``ricorri su $n$; al caso $\zero$ rispondi $m$; al caso $\suc\,k$ prendi il \emph{risultato gi\`a calcolato} su $k$ --- che chiamo $\mathit{rec}$, cio\`e $k + m$ --- e ne fai $\suc$''. Il parametro $\mathit{rec}$ \`e \textbf{esattamente} l'ipotesi induttiva $C\,k$; il suo essere disponibile \`e ci\`o che distingue $\indN$ dal puro case-split.

In Agda, con pattern matching, lo stesso si scrive nella forma familiare:
\begin{agda}
_+_ : ℕ → ℕ → ℕ
zero  + m = m
suc n + m = suc (n + m)   -- il "(n + m)" qui dentro È l'ipotesi induttiva,
                          --   cioè la chiamata ricorsiva su n
\end{agda}
Le due scritture sono equivalenti: la seconda \`e ci\`o che Agda elabora nella prima.

\section*{Case-split puro \emph{vs} induzione}

Per esplicitare ulteriormente la differenza con $\mathbb{B}$, ecco come sarebbe $\mathbb{N}$ con \emph{solo} case-split (l'analogo di $\indB$):
\[
\inferrule{
n : \mathbb{N} \\
C : \mathbb{N} \to \Type \\
b : C\,\zero \\
s : (k : \mathbb{N}) \to C\,(\suc\,k) \quad \textit{NB: $s$ non riceve $C\,k$}
}{
\mathsf{case\text{-}\mathbb{N}}(n, C, b, s) : C\,n
}
\]

Con questo non si dimostra quasi niente: nel ramo \texttt{suc k} non c'\`e ipotesi su \texttt{k}. La scelta di MLTT \`e di prendere l'\textbf{eliminatore induttivo} invece del solo case-split. Decisione di design, non un'inevitabilit\`a.

% ---------------------------------------------------------------------
\chapter{$\Pi$ --- la dipendenza}\label{cap:pi}

Le funzioni. Generalizzate al punto in cui il \emph{tipo di ritorno} pu\`o dipendere dal \emph{valore} in ingresso.

\section*{F --- Formazione}
\[
\inferrule{A : \Type \\ x : A \vdash B\,x : \Type}{(x : A) \to B\,x : \Type}
\]

$x : A \vdash B\,x : \Type$ significa: \emph{nel contesto in cui ho un $x$ di tipo $A$, $B\,x$ \`e un tipo}. $x$ \`e libera fuori; la regola la chiude dentro il tipo (Cero~\ref{cero:contesti} al lavoro).

Qui c'\`e la \textbf{decisione fondativa} di tutta MLTT:

\begin{cero}[Un termine pu\`o apparire dentro un tipo]\label{cero:dipendenza}
Questa \`e \textsc{la} decisione che definisce MLTT. Prima, hai i tipi semplici, dove $B$ non sa niente di $x$. Dopo, hai i tipi dipendenti.

Non \`e giustificata da niente di precedente. \`E una scelta su \emph{cosa vuoi che il sistema sia}. Potresti scegliere diversamente --- tipi semplici, sistema F, tipi lineari --- e avresti un sistema \textbf{diverso, non sbagliato}. Questo cero distingue MLTT da tutto il resto.
\end{cero}

\section*{I --- Introduzione}
\[
\inferrule{x : A \vdash b\,x : B\,x}{\lambda x.\,b\,x : (x : A) \to B\,x}
\]

Sotto l'ipotesi $x : A$ si costruisce un $b\,x : B\,x$; la regola scarica l'ipotesi producendo la lambda.

\section*{E --- Eliminazione}
\[
\inferrule{f : (x : A) \to B\,x \\ a : A}{f\,a : B\,a}
\]

L'applicazione. \textbf{Banale, nessuna ricorsione} --- $\Pi$ non \`e induttivo, non c'\`e struttura da srotolare in profondit\`a. Contrasto netto con $\indN$.

\section*{C --- Computazione}
\[
(\lambda x.\,b\,x)\,a \definitorio b\,a
\]

La $\beta$-riduzione. Definitoria (Cero~\ref{cero:comp}).

\section*{Nomi alternativi}
\begin{itemize}
\item \textbf{$\Pi$ type} --- nome formale MLTT
\item \textbf{dependent function type} --- descrittivo
\item \textbf{dependent product} --- nome storico, confonde
\end{itemize}

% ---------------------------------------------------------------------
\chapter{$\Pi$ in azione --- i tipi sono termini, ma in due passate}\label{cap:pi-azione}

Sezione di chiarimento sull'uso di $\Pi$. Tre cose vanno fissate, perch\'e Pie (\emph{The Little Typer}) e Agda le fondono sintatticamente.

\section*{I tipi sono termini --- vero, non illusione}

In MLTT esiste un universo $U$ (\agdainl{Set} in Agda) i cui abitanti sono tipi. Quindi una variabile $A : U$ \`e un termine \emph{onesto}: si passa come si passa un numero. \`E il senso vero di ``tipi di prima classe''.

\section*{Ma c'\`e una doppia passata}

Da \emph{The Little Typer}:
\begin{verbatim}
(claim elim-Pair
  (Pi ((A U) (D U) (X U))
      (-> (Pair A D) (-> A D X) X)))

(define elim-Pair
  (lambda (A D X)
    (lambda (p f) (f (car p) (cdr p)))))
\end{verbatim}

Quando il typechecker incontra \verb|(elim-Pair Nat Atom Bool p f)| fa \textbf{due lavori distinti}:

\begin{enumerate}
\item \textbf{Passata 1 --- type checking} (``a compile time''). Contano $A, D, X$. Il checker istanzia $A := \mathrm{Nat}$ ecc.\ nel claim per calcolare di che tipo devono essere \texttt{p} e \texttt{f}.
\item \textbf{Passata 2 --- computazione} (``a runtime/normalizzazione''). Contano \texttt{p} e \texttt{f}. Il corpo \verb|(f (car p) (cdr p))| srotola la coppia. Qui $A, D, X$ sono \textbf{gi\`a stati consumati} nella passata 1 --- non fanno pi\`u nulla.
\end{enumerate}

\textbf{Prima fissi i tipi, poi fai il resto.} L'illusione \emph{non} \`e che i tipi fingano di essere valori --- lo sono. L'illusione \`e che una sola sintassi \verb|(lambda (A D X) (lambda (p f) ...))| nasconde che le due $\lambda$ giocano ruoli temporalmente diversi.

\begin{oss}[Segnale fisico]
$A, D, X$ \emph{non compaiono nel corpo}. Un argomento che il calcolo non usa, ma che il checker pretende, \`e per definizione della passata-tipi.
\end{oss}

\begin{oss}[Conferma teorica --- \emph{type erasure}]
Proprio perch\'e la passata 2 non guarda $A, D, X$, un compilatore pu\`o cancellarli dopo il type checking. Se le passate fossero la stessa cosa, non si potrebbe. Il fatto che si possa \`e la prova che sono separate. (Dove le passate si \emph{toccano} \`e \texttt{Vec A n} (\S\ref{cap:vec}): l\`i un valore \texttt{n} risale dentro un tipo. Si intrecciano, non collassano.)
\end{oss}

\section*{\texttt{elim-Pair} in Agda}

Il claim Pie diventa:

\begin{agda}
-- Pair A D = Sigma A (\ _ -> D)  (prodotto NON dipendente; vedi cap. sigma)

-- versione esplicita: traduzione fedele del claim Pie
elim-Pair : (A : Set) (D : Set) (X : Set) → Pair A D → (A → D → X) → X
elim-Pair A D X p f = f (fst p) (snd p)
--        ^^^^^                          passata 1: tipi, NON usati nel corpo
--              ^^^                       passata 2: valori, USATI nel corpo

-- versione implicita: Agda inferisce A, D, X dal tipo di p
elim-Pair : {A D X : Set} → Pair A D → (A → D → X) → X
elim-Pair p f = f (fst p) (snd p)
\end{agda}

Le graffe \texttt{\{\}} \emph{sono} il riconoscimento sintattico che $A, D, X$ vivono nella passata 1. Chiamando \texttt{elim-Pair miaCoppia miaFunzione}, Agda guarda il tipo di \texttt{miaCoppia} (es. \texttt{Pair Nat Atom}), unifica e riempie $A := \mathrm{Nat}$, $D := \mathrm{Atom}$ da s\'e. La passata 1 avviene, ma invisibile.

\section*{$\Pi$ \`e unario --- gli argomenti multipli sono curry}

Dubbio naturale: \texttt{(A : Set)(D : Set)(X : Set) ->} sembra un $\Pi$ a tre argomenti. \textbf{Non lo \`e.} La regola di formazione di \S\ref{cap:pi} resta \textbf{unaria}. Quella scrittura \`e zucchero per $\Pi$ annidati:

\begin{agda}
(A : Set) → ((D : Set) → ((X : Set) → (Pair A D → ((A → D → X) → X))))
\end{agda}

Un solo $\Pi$ per ogni \texttt{->}, ciascuno con il contesto esteso di un'ipotesi (Cero~\ref{cero:contesti}).

\begin{oss}[Prova che \`e curry e non un $\Pi$ ternario primitivo]
Le applicazioni parziali hanno un tipo.
\begin{agda}
elim-Pair Nat            -- : (D : Set) (X : Set) → Pair Nat D → (Nat → D → X) → X
elim-Pair Nat Atom       -- : (X : Set) → Pair Nat Atom → (Nat → Atom → X) → X
elim-Pair Nat Atom Bool  -- : Pair Nat Atom → (Nat → Atom → Bool) → Bool
\end{agda}

Se fosse un $\Pi$ ternario primitivo, \texttt{elim-Pair Nat} da solo non avrebbe tipo. Il fatto che ce l'ha \`e la dimostrazione che \`e curry --- \textbf{curry dipendente}, perch\'e ogni $\Pi$ vede le variabili legate dai precedenti.
\end{oss}

\bigskip

\noindent\textbf{Da fissare.} \texttt{(A : ..)(B : ..)(C : ..) -> ...} = tre $\Pi$ annidati, non una regola nuova.

% ---------------------------------------------------------------------
\chapter{Il $\forall$ di Agda}\label{cap:forall}

Solo zucchero in pi\`u, ma \`e bene fissarlo perch\'e lo si vede dappertutto.

\begin{agda}
foo : ∀ (x : A) → B x       -- identico a (x : A) → B x
foo : ∀ x → B x             -- identico, ma il tipo di x è inferito da B
\end{agda}

% MODIFICATO: esplicitato che $\forall$ NON è un quantificatore logico primitivo
$\forall$ \textbf{non \`e un quantificatore logico primitivo, n\'e una regola distinta da $\Pi$}: \`e \emph{soltanto zucchero sintattico} per $\Pi$ con il permesso di omettere annotazioni di tipo (quando Agda pu\`o inferirle). Il nome viene dalla lettura logica (``per ogni $x$\ldots''), ma in MLTT non esiste un ``per ogni'' di prima classe separato da $\Pi$: la quantificazione universale \emph{\`e} $\Pi$, via Curry--Howard. Si applica sempre la regola di \S\ref{cap:pi}. Nessuna regola nuova.

% ---------------------------------------------------------------------
\chapter{$\Sigma$ --- il duale di $\Pi$}\label{cap:sigma}

Coppie dipendenti: il secondo componente ha un tipo che dipende dal primo.

\section*{F --- Formazione}
\[
\inferrule{A : \Type \\ x : A \vdash B\,x : \Type}{\Sigma\,x \in A,\,B\,x : \Type}
\]

Stessa premessa di $\Pi$ --- eredita il Cero~\ref{cero:dipendenza}, non ne aggiunge.

\section*{I --- Introduzione}
\[
\inferrule{a : A \\ b : B\,a}{(a, b) : \Sigma\,x \in A,\,B\,x}
\]

\section*{E --- Eliminazione}
\[
\inferrule{
p : \Sigma\,x \in A,\,B\,x \\
C : (\Sigma\,x \in A,\,B\,x) \to \Type \\
f : (a : A) \to (b : B\,a) \to C\,(a, b)
}{
\indS(p, C, f) : C\,p
}
\]

\textbf{Nessuna ricorsione} nel lato destro: una coppia ha esattamente due elementi, non c'\`e struttura da srotolare in profondit\`a. Come $\Pi$, contrasto con $\indN$.

\section*{C --- Computazione}
\[
\indS((a, b), C, f) \definitorio f\,a\,b
\]

\section*{Caso non dipendente: il prodotto}

Quando $B$ non usa $x$, si ottiene il prodotto ordinario:
\[
A \times B \;\;\coloneqq\;\; \Sigma\,A\,(\lambda\,\_\,.\,B)
\]

\texttt{Pair A D} di \S\ref{cap:pi-azione} era questo. Gli accessori \texttt{fst} e \texttt{snd} sono casi particolari di $\indS$ con motivo costante.

\section*{La confusione prodotto/somma}

\begin{center}
\renewcommand{\arraystretch}{1.15}
\begin{tabular}{lll}
\toprule
\textit{Senso} & \textit{Prodotto} & \textit{Somma} \\
\midrule
aritmetico/insiemistico & coppie $A \times B$ & scelta $A + B$ \\
categoriale             & $\Pi$ (funzioni indicizzate) & $\Sigma$ (coppie indicizzate) \\
\bottomrule
\end{tabular}
\end{center}

% MODIFICATO: la confusione "Σ si chiama somma in CT" spostata in nota a pie pagina
\noindent $\Sigma$ generalizza il prodotto \emph{aritmetico} ma si chiama \emph{somma categoriale}\footnote{In teoria delle categorie ``somma'' (o \emph{coprodotto}) \`e il duale di ``prodotto''. Il coprodotto binario \`e $A + B$, e la sua versione indicizzata \`e $\coprod_{x \in A} B\,x$, che \emph{\`e} $\Sigma\,x\!\in\!A,\,B\,x$. Quindi $\Sigma$ \`e ``somma'' nel senso categoriale di coprodotto indicizzato, anche se insiemisticamente i suoi abitanti sono coppie e ``sembra'' un prodotto. Doppia tradizione terminologica, da non confondere.}. Incidente storico, da ignorare. In pratica: $\Pi$ = funzioni dipendenti, $\Sigma$ = coppie dipendenti.

% ---------------------------------------------------------------------
\chapter{$\_\identita\_$ --- il tipo identit\`a}\label{cap:identita}

\epigraph{Una volta capita questa sezione, MLTT smette di sembrare misterioso.}{}

\noindent {\spacedlowsmallcaps{Cardine}}. Qui si chiude il cerchio: anche l'uguaglianza ha le sue quattro regole, \`e un tipo come un altro, e si distingue nettamente dall'uguaglianza definitoria del metalinguaggio.

\bigskip

Fino a \S\ref{cap:giudizi} e \S\ref{cap:test} abbiamo parlato di un'uguaglianza che vive \emph{nel metalinguaggio}: scritta $\equiv$, \`e un giudizio, non un tipo, non si abita, non si dimostra. \`E quella che dice ``queste due cose riducono alla stessa cosa per le regole di computazione''. La chiamiamo \textbf{uguaglianza definitoria}.

Ora introduciamo un \emph{secondo} $\equiv$, completamente diverso, che vive \textbf{nel linguaggio}: \`e un tipo, lo si abita con prove, lo si manipola. Si chiama \textbf{tipo identit\`a} (o \textbf{uguaglianza proposizionale}). Si scrive uguale e questa \`e una sciagura notazionale storica, ma sono due bestie diverse.

\medskip

Per smarcarli scrivo:
\begin{itemize}
\item $a \definitorio b : A$ per l'uguaglianza \textbf{definitoria} (metalinguaggio, \S\ref{cap:giudizi});
\item $a \identita b$ (senza il ``$: \Type$''), o $\mathsf{Id}\,A\,a\,b$ se serve esplicito, per il \textbf{tipo identit\`a} (linguaggio, questa sezione).
\end{itemize}

Quando Agda scrive \agdainl{a ≡ b} intende \textbf{il tipo identit\`a}, sempre. La cosa che gli viene dietro \`e il \emph{secondo} dei due.

\section*{Le quattro regole}

\paragraph{F --- Formazione.}
\[
\inferrule{A : \Type \\ x : A \\ y : A}{x \identita y : \Type}
\]
Dati due elementi di \emph{uno stesso} tipo, si ottiene un tipo: il tipo delle prove che sono uguali.

\paragraph{I --- Introduzione.} (Un solo costruttore --- e questo \`e il punto.)
\[
\inferrule{A : \Type \\ x : A}{\refl : x \identita x}
\]

C'\`e \textbf{un solo} costruttore, $\refl$, e abita \textbf{soltanto} $x \identita x$. Cio\`e: l'unico modo canonico di costruire una prova di uguaglianza \`e quando i due lati sono \textbf{gi\`a la stessa cosa}.

\bigskip

\noindent\textbf{Qui salta fuori la cosa cruciale}: ``gi\`a la stessa cosa'' \emph{in che senso}? Nel senso dell'\textbf{uguaglianza definitoria del metalinguaggio}, cio\`e dopo aver applicato tutte le regole di computazione viste finora ($\indN$ su \texttt{zero}, $\beta$-riduzione di $\Pi$, ecc.).

\begin{regola}[Il ponte]\label{reg:ponte}
$\refl$ ha tipo $x \identita y$ (linguaggio) \textbf{se e solo se} $x \definitorio y : A$ (definitorio, metalinguaggio).
\end{regola}

$\refl$ \emph{colma la trincea} tra meta e linguaggio. \`E il termine che dice ``i due lati riducono allo stesso normale, e io ti porto questa informazione dentro il sistema sotto forma di prova''.

Quindi le due uguaglianze non sono indipendenti: il tipo identit\`a \emph{si appoggia} alla definitoria. La definitoria sta ``sotto'' e fa il lavoro silenzioso di riduzione; il tipo identit\`a sta ``sopra'' e ne raccoglie il risultato come prova manipolabile.

% NUOVO: esempi Agda di refl che funziona e refl che fallisce
\bigskip

\noindent\textbf{In pratica, con $+$ definita come in \S\ref{cap:nat}}:
\begin{agda}
-- (1) refl funziona: entrambi i lati riducono a 4
_ : 2 + 2 ≡ 4
_ = refl

-- (2) refl funziona ANCHE qui: 2+2 e 3+1 hanno lo stesso normale (= 4)
_ : 2 + 2 ≡ 3 + 1
_ = refl

-- (3) refl funziona: zero + n riduce a n per la prima equazione di +
_ : ∀ {n} → zero + n ≡ n
_ = refl

-- (4) refl FALLISCE: n + zero NON riduce, perché + ricorre sul PRIMO
--     argomento e n è una variabile (nessuna delle equazioni si attiva).
--     Definitoriamente non è uguale a n; serve induzione su n per dimostrarlo.
-- _ : ∀ {n} → n + zero ≡ n
-- _ = refl   -- errore di typecheck: n + zero != n
\end{agda}

\noindent (1)--(3) sono ``stessa normale dopo riduzione'', quindi la Regola~\ref{reg:ponte} li accetta. (4) non lo \`e: \emph{semanticamente} $n + \zero$ vale $n$, ma \emph{definitoriamente} non riduce, perch\'e la regola pertinente di $+$ \`e bloccata sulla variabile $n$. Serve $\indN$ per dimostrarlo --- e il risultato non sar\`a $\refl$, ma un termine costruito induttivamente.

\medskip

\noindent\textbf{Da fissare.} ``$\refl$ funziona'' $\iff$ ``i due lati riducono alla stessa forma normale per le regole di computazione''. Tutto il resto richiede prove genuine.

\paragraph{E --- Eliminazione.}
\[
\inferrule{
p : x \identita y \\
C : (y : A) \to (x \identita y) \to \Type \\
c : C\,x\,\refl
}{
\Jelim(p, C, c) : C\,y\,p
}
\]

L'eliminatore $\Jelim$ (\emph{path induction}). Per costruire qualcosa che dipende da una prova generica $p : x \identita y$, basta saperlo costruire nel caso canonico $p = \refl$ (dove necessariamente $y = x$).

\paragraph{C --- Computazione.}
\[
\Jelim(\refl, C, c) \definitorio c
\]
Quando la prova \`e canonica, $\Jelim$ restituisce il caso base.

% NUOVO: subst come corollario immediato di J --- J in azione
\subsection*{$\Jelim$ in azione --- $\mathrm{subst}$ come corollario immediato}

Il principio della \emph{path induction} \`e formidabile ma astratto. Il suo primo corollario \`e una funzione che usiamo continuamente: \textbf{trasporto} di una propriet\`a lungo un'uguaglianza.
\[
\mathrm{subst} : \{A : \Type\}\,(P : A \to \Type)\,\{x\,y : A\} \to x \identita y \to P\,x \to P\,y
\]
Letta a voce: ``se $x \identita y$ e ho un abitante di $P\,x$, ne ottengo uno di $P\,y$''. Si ottiene da $\Jelim$ scegliendo come motivo $C\,y\,p \definitorio P\,y$ (cio\`e: il motivo dipende solo dall'estremo, ignora la prova):
\begin{agda}
subst : {A : Set} (P : A → Set) {x y : A} → x ≡ y → P x → P y
subst P p px = J p (λ y _ → P y) px
--                  ^^^^^^^^^^^
--                  motivo: "non mi importa della prova,
--                   voglio dire qualcosa su y"
\end{agda}
\noindent Il caso base $C\,x\,\refl \definitorio P\,x$ \`e gi\`a dato (lo \`e \texttt{px}); la computazione di $\Jelim$ su $\refl$ restituisce esattamente quel \texttt{px}. Quando invece la prova non \`e $\refl$ ma una qualunque, $\Jelim$ ``trasporta'' il dato lungo l'identit\`a.

\begin{oss}[Perch\'e questo importa]
$\mathrm{subst}$ \`e ci\`o che permette il \emph{rewriting} nelle dimostrazioni in Agda (la parola chiave \texttt{rewrite} si elabora in chiamate a $\mathrm{subst}$). Senza $\Jelim$, non ci sarebbe nessun modo di ``usare'' una prova di uguaglianza per cambiare il tipo di un termine. $\Jelim$ \`e l'astratto; $\mathrm{subst}$ \`e il primo concreto.
\end{oss}

\section*{Ricapitolazione della distinzione}

\begin{center}
\renewcommand{\arraystretch}{1.2}
\begin{tabular}{lll}
\toprule
 & \textit{Uguaglianza definitoria} & \textit{Tipo identit\`a} \\
\midrule
Dove vive & \textbf{metalinguaggio} & \textbf{linguaggio} \\
Forma     & giudizio $a \definitorio b : A$ & tipo $a \identita b$ \\
Si abita? & no, \`e un'affermazione & s\`i, con $\refl$ (e via $\Jelim$) \\
Si scrive? & nelle regole, non in Agda & in Agda: \agdainl{a ≡ b} \\
Viene da  & regole di computazione (Cero~\ref{cero:comp}) & regole F-I-E-C di \S\ref{cap:identita} \\
Ponte     & \multicolumn{2}{c}{abitato da $\refl$ $\iff$ vale la definitoria (Regola~\ref{reg:ponte})} \\
\bottomrule
\end{tabular}
\end{center}

\bigskip

\noindent Quindi, quando in Agda si scrive:
\begin{agda}
qualcosa : x ≡ y
qualcosa = refl
\end{agda}
si \emph{abita} un tipo (linguaggio); Agda accetta solo se \texttt{x} e \texttt{y} sono definitoriamente uguali (metalinguaggio). I due livelli parlano tramite $\refl$. \textbf{Questa \`e la cosa pi\`u importante da tenere chiara di tutto MLTT.}

\section*{Quarto punto dove si inciampa}

In \S\ref{cap:test} avevamo promesso un quarto punto, da chiudere appena introdotto il tipo identit\`a. Eccolo: i due $\equiv$ si scrivono uguale ma sono due cose diverse. Sopra c'\`e la regola per distinguerli --- guarda \emph{cosa} stai dicendo: un'affermazione sul sistema, o un tipo che vuoi abitare?

\medskip

Completiamo la tabella lampo di \S\ref{cap:test}:

\begin{center}
\renewcommand{\arraystretch}{1.15}
\begin{tabular}{lll}
\toprule
\textit{Vedi questo} & \textit{Dov'\`e} & \textit{Perch\'e} \\
\midrule
$\equiv$ in $(\lambda x.\,b)\,a \equiv b\,a$ (regola di computazione) & meta & uguaglianza definitoria, asserzione \\
$\equiv$ in $x \equiv y : \Type$ (Agda: \texttt{\_$\equiv$\_}) & linguaggio & tipo identit\`a, si abita con $\refl$ \\
$\refl$ & linguaggio & termine, abita il tipo identit\`a \\
$\indB$, $\indN$, $\indS$, $\Jelim$ & meta & eliminatori = regole di inferenza \\
\bottomrule
\end{tabular}
\end{center}

% ---------------------------------------------------------------------
\chapter{$\mathrm{Vec}\,A\,n$ --- tutto messo insieme}\label{cap:vec}

Capstone della Parte~II: primo tipo dove $\mathbb{N}$, $\Pi$ e la dipendenza Cero~\ref{cero:dipendenza} si incastrano davvero. \texttt{Vec A n} \`e il tipo delle liste di elementi di $A$ con \emph{esattamente} $n$ elementi. \textbf{$n$ vive nel tipo} --- il typechecker conosce la lunghezza a compile time. Questo \`e il Cero~\ref{cero:dipendenza} al lavoro su un valore numerico (\`e anche il punto dove le due passate di \S\ref{cap:pi-azione} si \emph{toccano}: $n$ \`e un valore che entra in un tipo).

\section*{F --- Formazione}
\[
\inferrule{A : \Type \\ n : \mathbb{N}}{\mathrm{Vec}\,A\,n : \Type}
\]

Il tipo dipende da un \emph{valore} $n : \mathbb{N}$. Non \`e \texttt{Vec A} genericamente: \`e un tipo diverso per ogni lunghezza.

\section*{I --- Introduzione}
\[
\inferrule{ }{\nil : \mathrm{Vec}\,A\,\zero}
\qquad
\inferrule{a : A \\ as : \mathrm{Vec}\,A\,n}{\cons\,a\,as : \mathrm{Vec}\,A\,(\suc\,n)}
\]

\texttt{cons} incrementa $n$ nel tipo risultante. Non per calcolo a runtime, \textbf{per costruzione}.

\section*{E --- Eliminazione (metalinguaggio)}
\[
\inferrule{
vs : \mathrm{Vec}\,A\,n \\
C : (n : \mathbb{N}) \to \mathrm{Vec}\,A\,n \to \Type \\
b : C\,\zero\,\nil \\
s : (k : \mathbb{N}) \to (a : A) \to (as : \mathrm{Vec}\,A\,k) \to C\,k\,as \to C\,(\suc\,k)\,(\cons\,a\,as)
}{
\indV(vs, C, b, s) : C\,n\,vs
}
\]

Stessa struttura di $\indN$: caso base, e passo che riceve l'ipotesi induttiva $C\,k\,as$. Il motivo $C$ dipende sia dalla lunghezza che dal vettore.

\section*{C --- Computazione}
\begin{align*}
\indV(\nil,         C, b, s) &\definitorio b \\
\indV(\cons\,a\,as, C, b, s) &\definitorio s\,k\,a\,as\,(\indV(as, C, b, s))
\end{align*}

Ricorsione su \texttt{as}, strutturalmente pi\`u corto di \texttt{cons a as}. Stesso schema di $\indN$.

\section*{Nota sintattica: parametri \emph{vs} indici}

\begin{agda}
data Vec (A : Set) : ℕ → Set where
--       ^^^^^^^^   ^^^^^^^^^
--       parametro  indice
\end{agda}

\begin{itemize}
\item \textbf{Parametro} \texttt{(A : Set)}: prima dei \texttt{:}, fisso per tutti i costruttori.
\item \textbf{Indice} \texttt{: $\mathbb{N}$ -> Set}: dopo i \texttt{:}, pu\`o cambiare costruttore per costruttore (\texttt{nil} lo fissa a \texttt{zero}, \texttt{cons} lo porta da $n$ a $\suc\,n$).
\end{itemize}

Se $n$ fosse un parametro, \texttt{cons} non potrebbe cambiarlo. Questa \`e una decisione di design (\S\ref{cap:ceri}).

% NUOVO: avvertimento sull'inferenza fragile per i tipi con indici
\section*{Avvertimento --- inferenza fragile sugli indici}

L'inferenza di tipo di Agda \`e completa per i \emph{parametri}, ma \textbf{non lo \`e per gli indici}. Quando un'operazione richiede che due indici risultino definitoriamente uguali e non lo sono per pura riduzione, l'unificatore si arresta e devi fornire una prova di uguaglianza (tipicamente con \texttt{rewrite} o $\mathrm{subst}$, \S\ref{cap:identita}). Tre situazioni tipiche:

\begin{enumerate}
\item \textbf{Concatenazione e lemmi di associativit\`a/commutativit\`a sugli indici}:
\begin{agda}
_++_ : ∀ {A m n} → Vec A m → Vec A n → Vec A (m + n)

-- e ora, per dimostrare l'associatività di ++ a livello di vettori,
-- serve PRIMA un lemma su ℕ (associatività di +) e POI un rewrite,
-- perché (m + n) + p non è DEFINITORIAMENTE uguale a m + (n + p):
-- lo è solo proposizionalmente, e Agda non può forzarlo da solo.
\end{agda}

\item \textbf{Funzioni con indice non-canonico in posizione di pattern}:
\begin{agda}
head : ∀ {A n} → Vec A (suc n) → A
head (cons a _) = a
-- ok: l'indice è "suc n", l'unificazione vede subito che il vettore
-- non può essere nil (che vorrebbe l'indice zero). Bene.

-- Ma se chiami head su un Vec A m con m generico, prima devi
-- convincere Agda che m ≡ suc k per qualche k.
\end{agda}

\item \textbf{Funzioni la cui forma del risultato non si decide per riduzione}: capita spesso quando l'indice contiene una funzione bloccata su una variabile (come $n + \zero$ rispetto a $n$).
\end{enumerate}

\noindent\textbf{Da fissare.} Tipi dipendenti su indici sono potenti, ma il prezzo \`e che l'inferenza non \`e ``magica'': dove la riduzione si ferma, devi fornire tu il ponte con una prova proposizionale. Questo non \`e un difetto di Agda; \`e la conseguenza diretta della Regola~\ref{reg:ponte} --- definitorio sta sotto, proposizionale sopra.

% NUOVO: tabella riassuntiva finale di tutta la Parte II
\section*{Quadro d'insieme --- la Parte~II in una tabella}

I sei tipi della Parte~II, allineati sulle stesse colonne. Stessa anatomia F-I-E-C; ci\`o che cambia \`e la presenza (o l'assenza) di ricorsione, e da cosa dipende il motivo $C$.

\begin{center}
\renewcommand{\arraystretch}{1.25}
\small
\begin{tabular}{llllp{4.6cm}}
\toprule
\textit{Tipo} & \textit{F-I-E-C} & \textit{Ricorsivo} & \textit{Motivo dipende da} & \textit{Uso principale} \\
\midrule
$\mathbb{B}$ & s\`i & no & $b : \mathbb{B}$ & case-split puro; pi\`u piccolo induttivo possibile \\
$\mathbb{N}$ & s\`i & s\`i & $n : \mathbb{N}$ & induzione e ricorsione primitiva \\
$\Pi$        & s\`i & no & --- (non induttivo) & funzioni dipendenti; decisione fondativa MLTT \\
$\Sigma$     & s\`i & no & $p : \Sigma\,x\!\in\!A,\,B\,x$ & coppie dipendenti; duale di $\Pi$ \\
$\_\identita\_$ & s\`i & no & $y : A$ \emph{e} $p : x \identita y$ & uguaglianza interna; ponte meta--linguaggio \\
$\mathrm{Vec}\,A\,n$ & s\`i & s\`i & $n : \mathbb{N}$ \emph{e} $vs : \mathrm{Vec}\,A\,n$ & dipendenza su un valore numerico; capstone \\
\bottomrule
\end{tabular}
\end{center}

\medskip

\noindent\textbf{Letture trasversali.}
\begin{itemize}
\item \emph{Ricorsione = costruttore che cita il tipo stesso.} $\mathbb{B}$ e $\Sigma$ non hanno tale costruttore; $\mathbb{N}$ ($\suc$) e $\mathrm{Vec}$ ($\cons$) s\`i. $\Pi$ e $\_\identita\_$ non sono induttivi in senso stretto.
\item \emph{Motivo a una variabile vs.\ a due.} $\indB$, $\indN$, $\indS$ hanno motivi su una variabile (il valore eliminato). $\Jelim$ e $\indV$ ne hanno due, perch\'e il tipo dipende da pi\`u di un dato (in $\Jelim$: l'altro estremo e la prova; in $\mathrm{Vec}$: la lunghezza e il vettore).
\item \emph{Dipendenza Cero~\ref{cero:dipendenza}.} $\Pi$ la introduce; $\Sigma$ la eredita; $\_\identita\_$ e $\mathrm{Vec}$ la sfruttano in modo essenziale (un \emph{valore} entra in un \emph{tipo}). $\mathbb{B}$ e $\mathbb{N}$ no, di per s\'e: la usano solo quando si scelgono motivi non costanti.
\end{itemize}

\noindent Da qui in poi (Parte~III) tutti gli ingredienti sono pronti: le dimostrazioni saranno termini in tipi costruiti con queste sei regole F-I-E-C, niente di pi\`u.